% !TeX program = xelatex
% Απαιτεί XeLaTeX (λόγω fontspec/polyglossia) -- ΔΕΝ μεταγλωττίζεται με pdflatex.
\documentclass[11pt,a4paper]{article}

\usepackage{fontspec}
\usepackage{polyglossia}
\setmainlanguage{greek}
\setotherlanguage{english}
\setmainfont{Linux Libertine O}
\setsansfont{Linux Biolinum O}
\setmonofont{DejaVu Sans Mono}[Scale=0.9]

\usepackage[a4paper,margin=2.5cm]{geometry}
\usepackage{parskip}
\usepackage{amsmath,amssymb}
\usepackage{booktabs}
\usepackage{array}
\usepackage{longtable}
\usepackage{graphicx}
\graphicspath{{../results/}}
\usepackage{enumitem}
\usepackage{caption}
\usepackage[hidelinks,colorlinks=true,linkcolor=blue!50!black,citecolor=blue!50!black,urlcolor=blue!50!black]{hyperref}
\usepackage{xcolor}

\newcommand{\file}[1]{\texttt{#1}}
\newcommand{\var}[1]{\texttt{#1}}

\let\oldunderscore\_
\renewcommand{\_}{\oldunderscore\allowbreak}
\setlength{\emergencystretch}{3em}

\setlist[itemize]{topsep=3pt,itemsep=3pt}
\setlist[enumerate]{topsep=3pt,itemsep=3pt}

\title{\textbf{Online Convex Optimization based tuning of\\ Mixture of Experts for Sports Outcome Prediction}\\[8pt]
\large Αναφορά προόδου: τι έχει υλοποιηθεί}
\author{[Ονοματεπώνυμο φοιτητή]}
\date{\today}

\begin{document}
\maketitle

% =====================================================================
\section{Σύνοψη}
% =====================================================================

Η εργασία \textbf{δεν} εκπαιδεύει ένα ακόμη μοντέλο πρόβλεψης πάνω σε στατιστικά αγώνων ---
αυτό το κάνουν ήδη οι στοιχηματικές εταιρείες, με ασύγκριτα περισσότερα δεδομένα και
υπολογιστική ισχύ από μια διπλωματική. Αντ' αυτού αντιμετωπίζει \textbf{κάθε στοιχηματική
εταιρεία ως έναν \emph{expert}} που δίνει πιθανότητες $(p_H, p_D, p_A)$ για το αποτέλεσμα
κάθε αγώνα, και μαθαίνει \textbf{online} --- αγώνα προς αγώνα, χωρίς επανεκπαίδευση πάνω στο
παρελθόν --- πώς να τους συνδυάζει σε ένα μείγμα που τους ξεπερνά. Αυτό είναι το κλασικό
πλαίσιο \emph{prediction with expert advice} / \emph{Online Convex Optimization} (OCO), και
δίνει θεωρητικές εγγυήσεις regret χωρίς καμία στατιστική υπόθεση για την κατανομή των
δεδομένων.

Τι υπάρχει αυτή τη στιγμή:

\begin{itemize}
  \item Πλήρες pipeline που τρέχει end-to-end πάνω σε \textbf{76.584 αγώνες}, 22
        πρωταθλήματα, 10 σεζόν (2016/17--2025/26), \textbf{26 bookmakers} ως experts.
  \item Τρεις αλγόριθμοι OCO (\textbf{Hedge}, \textbf{OGD}, \textbf{FTRL}), γενικευμένοι
        ώστε να δουλεύουν με bookmakers μερικής κάλυψης (αναγωγή \emph{Sleeping Experts}).
  \item Δεύτερο στάδιο OCO για \textbf{βαθμονόμηση} (online temperature scaling), και τρίτο
        για το \textbf{μέγεθος στοιχήματος} (online Kelly).
  \item \textbf{Στατιστικός έλεγχος} με moving block bootstrap σε κάθε σύγκριση --- κανένα
        συμπέρασμα δεν στηρίζεται σε «φαίνεται καλύτερο».
  \item \textbf{9 επεκτάσεις} και \textbf{4 έλεγχοι μεθοδολογικής ευρωστίας}, ο καθένας
        απαντά ένα συγκεκριμένο ερώτημα (κατάλογος στον Πίνακα~\ref{tab:built}).
\end{itemize}

Τα τρία κύρια ευρήματα, σε μία γραμμή το καθένα:

\begin{enumerate}
  \item Το μείγμα \textbf{ξεπερνά στατιστικά σημαντικά κάθε μεμονωμένο bookmaker εκτός
        ενός}: του PSC (Pinnacle, closing odds). Στη διαμόρφωση που έχει πλέον επιλεγεί ως
        καλύτερη --- panel μόνο από closing αποδόσεις --- \textbf{η διαφορά από τον PSC
        παύει να είναι στατιστικά σημαντική} ($p=0{,}234$), δηλαδή το μείγμα τον φτάνει
        χωρίς να τον ξεπερνά αποδεδειγμένα.
  \item Η \textbf{βαθμονόμηση βοηθάει σημαντικά} και στους 4 αλγορίθμους --- το «δεύτερο
        στάδιο OCO» δεν είναι διακοσμητικό.
  \item Καλύτερο log-loss \textbf{δεν} σημαίνει κέρδος: η οικονομική προσομοίωση δίνει
        σημαντική άκρη σε 3 από 11 bookmakers και τίποτα στους υπόλοιπους 8 --- η άκρη
        υπάρχει απέναντι σε συγκεκριμένους bookmakers, όχι ως γενική στρατηγική.
\end{enumerate}

% =====================================================================
\section{Δεδομένα}
% =====================================================================

Πηγή: football-data.co.uk, ενότητα «Main Leagues» (η ενότητα «Extra Leagues» έχει άλλη μορφή
και συνήθως δεν περιέχει αποδόσεις, οπότε είναι εκτός σκοπού).

\begin{itemize}
  \item \textbf{22 πρωταθλήματα}: Αγγλία $\times$5, Σκωτία $\times$4,
        Γερμανία/Ιταλία/Ισπανία/Γαλλία $\times$2, Ολλανδία/Βέλγιο/Πορτογαλία/Τουρκία/Ελλάδα
        $\times$1.
  \item \textbf{10 σεζόν} (2016/17--2025/26), \textbf{76.639 αγώνες}, από τους οποίους
        \textbf{76.584} έχουν τουλάχιστον μία χρήσιμη απόδοση --- αυτοί είναι οι «γύροι» πάνω
        στους οποίους τρέχουν όλοι οι αλγόριθμοι.
  \item \textbf{26 ανεξάρτητοι bookmakers} μετά τον καθαρισμό: αποκλείονται οι
        συγκεντρωτικές στήλες (Max/Avg, Betbrain) ως μη ανεξάρτητοι experts, και
        συγχωνεύεται ένα γνωστό rebrand (VC Bet $\to$ BetVictor) σε μία συνεχή ταυτότητα.
  \item Η κάλυψη ανά bookmaker κυμαίνεται από \textbf{6,9\%} (CLC) έως \textbf{99,8\%}
        (B365) --- εξ ου και η ανάγκη για την αναγωγή Sleeping Experts. Συγκρίσεις
        θεωρούνται αξιόπιστες μόνο για κάλυψη $\geq$70\% («Tier A»), και οι bookmakers
        ενεργοί σε μία μόνο σεζόν επισημαίνονται ρητά ως μη συγκρίσιμοι.
\end{itemize}

Το panel χωρίζεται καθαρά σε \textbf{13 opening} και \textbf{13 closing} αποδόσεις (η
σύμβαση του football-data.co.uk: κατάληξη «C» σημαίνει την απόδοση κλεισίματος του ίδιου
bookmaker) --- διάκριση που αποδείχθηκε ουσιώδης, βλ.\ \S\ref{sec:robust}.

% =====================================================================
\section{Μεθοδολογικός σκελετός}
% =====================================================================

\paragraph{De-vig.} Οι δημοσιευμένες αποδόσεις περιέχουν το περιθώριο κέρδους του bookmaker
(overround): τα αντίστροφά τους αθροίζουν σε $>1$. Μετατρέπονται σε πιθανότητες με
proportional normalization, $p_i = (1/o_i) / \sum_j (1/o_j)$. Η επιλογή αυτή ελέγχθηκε
ανεξάρτητα με τη μέθοδο του Shin (1992), βλ.\ \S\ref{sec:robust}.

\paragraph{Πρώτο στάδιο OCO: τα βάρη του μείγματος.} Hedge (πολλαπλασιαστικά βάρη), OGD
(projected gradient descent στο simplex) και FTRL (L2 regularizer). Και οι τρεις δουλεύουν με
learning rate $\eta_t \propto 1/\sqrt{t}$, χωρίς κανένα hyperparameter προς συντονισμό σε
validation set.

\paragraph{Sleeping Experts / Specialists.} Οι κλασικοί αλγόριθμοι υποθέτουν ότι κάθε expert
δίνει πρόβλεψη σε κάθε γύρο, που εδώ δεν ισχύει (πολλοί bookmakers καλύπτουν $<$30\% των
αγώνων). Χρησιμοποιείται η κλασική αναγωγή (Freund, Schapire, Singer \& Warmuth, 1997): κάθε
expert κρατά μόνιμη κατάσταση, σε κάθε γύρο χρησιμοποιούνται και επανακανονικοποιούνται
\emph{μόνο} οι «ξύπνιοι», και η κατάσταση των «κοιμισμένων» παγώνει. Η εγγύηση regret
μετράται τότε μόνο στους γύρους όπου ο κάθε expert ήταν πραγματικά ενεργός.

\paragraph{Δεύτερο στάδιο OCO: βαθμονόμηση.} Διαγνωστικά βρέθηκε το γνωστό
\emph{favorite-longshot bias} σε \emph{όλες} τις σειρές, μείγματα συμπεριλαμβανομένων: οι
μικρές πιθανότητες είναι συστηματικά υπερεκτιμημένες και οι μεγάλες υποεκτιμημένες. Η
διόρθωση γίνεται με online temperature scaling, $\tilde p_k \propto p_k^{\,s}$: η συνάρτηση
απώλειας ως προς $s$ είναι log-sum-exp μιας γραμμικής συνάρτησης, άρα \emph{κυρτή}, οπότε το
$s$ μαθαίνεται με απλό projected OGD σε μία διάσταση --- ίδια συνταγή OCO, νέο στάδιο. Το
learning rate επιλέγεται \emph{αποκλειστικά} στο χρονολογικά πρώτο 75\% των γύρων, ώστε τα
αναφερόμενα νούμερα να μην έχουν δει ποτέ το validation κομμάτι.

\paragraph{Τρίτο στάδιο OCO (πείραμα): το μέγεθος στοιχήματος.} Αντί να μπαίνει η
εκτιμώμενη πιθανότητα στον κλειστό τύπο Kelly, το ποσοστό $f$ μαθαίνεται online με projected
OGD πάνω στην (κυρτή) απώλεια αρνητικού log-πλούτου --- η οπτική \emph{Online Portfolio
Selection} (Cover, 1991).

\paragraph{Στατιστική σημαντικότητα.} Οι διαφορές log-loss είναι της τάξης $10^{-4}$, οπότε
κάθε σύγκριση περνά από \textbf{moving block bootstrap} (2.000 δείγματα, μήκος block
$\approx n^{1/3}$). Ο λόγος που δεν χρησιμοποιείται απλό t-test ή i.i.d.\ bootstrap: οι
απώλειες διαδοχικών αγώνων \emph{δεν} είναι ανεξάρτητες (ένα αγωνιστικό διήμερο με πολλές
ανατροπές χτυπά ταυτόχρονα κάθε προγνώστη), οπότε η δειγματοληψία μεμονωμένων γύρων θα
υποεκτιμούσε τη διασπορά. Το block bootstrap δειγματοληπτεί συνεχόμενα μπλοκ και διατηρεί
αυτή τη βραχυχρόνια εξάρτηση --- η ίδια ιδέα με το κλασικό τεστ Diebold--Mariano, χωρίς
ασυμπτωτική υπόθεση κανονικότητας.

% =====================================================================
\section{Τι έχει υλοποιηθεί}
% =====================================================================

Το project είναι μια ακολουθία αυτόνομων σεναρίων Python που επικοινωνούν μέσω αρχείων CSV,
με ένα κοινό αρχείο-βιβλιοθήκη (\file{sleeping\_experts.py}) απ' όπου 15 άλλα εισάγουν
φόρτωση δεδομένων και αλγορίθμους. Κάθε επέκταση είναι ξεχωριστό αρχείο και απαντά
\emph{ένα} ερώτημα.

\begin{longtable}{@{}p{0.30\textwidth}p{0.64\textwidth}@{}}
\caption{Τι υπάρχει και τι ερώτημα απαντά.}\label{tab:built}\\
\toprule
\textbf{Συνιστώσα} & \textbf{Ερώτημα / περιεχόμενο} \\
\midrule
\endfirsthead
\toprule
\textbf{Συνιστώσα} & \textbf{Ερώτημα / περιεχόμενο} \\
\midrule
\endhead
\multicolumn{2}{@{}l}{\textbf{Πυρήνας}}\\[2pt]
Αγωγός δεδομένων & Κατέβασμα 220 αρχείων, de-vig, συγχώνευση rebrand, δύο εναλλακτικές
  κανονικοποιήσεις σε ένα πέρασμα. \\
Hedge / OGD / FTRL + Sleeping Experts & Ποιος αλγόριθμος OCO συνδυάζει καλύτερα τους
  experts, και πώς συγκρίνεται με τον καθένα τους χωριστά. \\
Ανάλυση βαθμονόμησης & Είναι οι πιθανότητες αξιόπιστες; (Διάγνωση: όχι --- favorite-longshot
  bias.) \\
Διόρθωση βαθμονόμησης & Το διορθώνει το online temperature scaling; \\
Τελική κατάταξη & Μία αυθεντική κατάταξη, raw και βαθμονομημένη, υπολογισμένες με τον
  \emph{ίδιο} τρόπο ώστε να συγκρίνονται δίκαια. \\
Έλεγχοι σημαντικότητας & Μπαταρία 12 συγκρίσεων με moving block bootstrap. \\[4pt]
\multicolumn{2}{@{}l}{\textbf{Επεκτάσεις}}\\[2pt]
Fixed-Share & Διορθώνει το «παγωμένο πλεονέκτημα»: σπάνιοι experts κρατούν διογκωμένο βάρος
  επειδή τίποτα δεν το διαβρώνει όσο κοιμούνται. \\
Contextual experts & Βοηθάει η εκμάθηση χωριστών βαρών ανά πρωτάθλημα; \\
Value betting & Θα έβγαζε λεφτά ένας πραγματικός bettor; (Kelly, leave-one-bookmaker-out.) \\
Online Kelly & Βοηθάει να μαθαίνεται και το \emph{ποσό} του στοιχήματος online; \\
Ακρίβεια ανά κατηγορία & Ποιος είναι καλός σε 1, ποιος σε Χ, ποιος σε 2 χωριστά. \\
Χρονικές τάσεις & Γίνεται η αγορά πιο αξιόπιστη με τα χρόνια; (Mann-Kendall.) \\
Παράθυρο pooling & Αλλάζουν τα συμπεράσματα με 5ετία αντί 10ετίας; \\
Βαθμονόμηση ανά σεζόν & Βοηθάει το reset της βαθμονόμησης κάθε σεζόν; \\
Warm-start βαθμονόμηση & Και αν μεταφέρεται η θερμοκρασία της προηγούμενης σεζόν; \\[4pt]
\multicolumn{2}{@{}l}{\textbf{Έλεγχοι ευρωστίας}}\\[2pt]
Shin έναντι proportional de-vig & Είναι τα ευρήματα artifact της μεθόδου de-vig; \\
Opening έναντι closing & Κάνει καλύτερο μείγμα η τελική τιμή της αγοράς από την αρχική; \\
Μεταφορά ποιότητας & Προβλέπει η «οξύτητα» στο closing την οξύτητα στο opening; \\
Έλεγχος look-ahead bias & Ήταν το σχέδιο leave-one-out πραγματικά αιτιακό; (Όχι σε μία
  περίπτωση --- διορθώθηκε, βλ.\ \S\ref{sec:robust}.) \\
\bottomrule
\end{longtable}

% =====================================================================
\section{Κύρια ευρήματα}
% =====================================================================

\subsection{Κατάταξη}

\begin{table}[h]
\centering
\begin{tabular}{@{}llrr@{}}
\toprule
\textbf{Σειρά} & \textbf{Τύπος} & \textbf{Raw log-loss} & \textbf{Βαθμονομημένο} \\
\midrule
PSC (Pinnacle, closing) & bookmaker & 0,99714 & \textbf{0,99685} \\
\textbf{OGD}            & αλγόριθμος & 0,99745 & \textbf{0,99693} \\
Hedge                   & αλγόριθμος & 0,99845 & 0,99776 \\
FTRL                    & αλγόριθμος & 0,99850 & 0,99780 \\
Uniform average         & baseline   & 0,99886 & 0,99813 \\
VC\_BV                  & bookmaker & 1,00075 & 1,00026 \\
WH                      & bookmaker & 1,00119 & 1,00048 \\
PS                      & bookmaker & 1,00083 & 1,00055 \\
IW                      & bookmaker & 1,00186 & 1,00061 \\
BW                      & bookmaker & 1,00140 & 1,00065 \\
B365                    & bookmaker & 1,00113 & 1,00073 \\
\bottomrule
\end{tabular}
\caption{Tier-A κατάταξη (κάλυψη $\geq$70\%), ταξινομημένη κατά βαθμονομημένο log-loss.
Χαμηλότερο = καλύτερο. Πηγή: \file{results/final\_ranking\_table.csv}.}
\end{table}

\begin{figure}[h]
\centering
\includegraphics[width=0.92\textwidth]{final_ranking.png}
\caption{Κάθε σειρά: κενός κύκλος = raw log-loss, γεμάτος = μετά τη βαθμονόμηση. Η επιλογή
dot/slope plot αντί ραβδογράμματος είναι σκόπιμη: τα log-loss μαζεύονται γύρω από το $1{,}0$
χωρίς ουσιαστικό μηδέν, οπότε ένα ραβδόγραμμα με κομμένο άξονα θα μεγαλοποιούσε οπτικά
μικροσκοπικές διαφορές.}
\end{figure}

\subsection{Τι είναι στατιστικά σημαντικό}

\begin{itemize}
  \item Και οι \textbf{τέσσερις} μέθοδοι μείγματος ξεπερνούν σημαντικά την ομοιόμορφη μέση
        τιμή (OGD: $-0{,}00141$· Hedge $-0{,}00041$· FTRL $-0{,}00037$, όλα $p<0{,}001$). Το
        \textbf{OGD} είναι ο καλύτερος των τριών, σημαντικά καλύτερο και από Hedge
        ($-0{,}00100$) και από FTRL ($-0{,}00104$), και ξεπερνά σημαντικά και τον B365C
        ($-0{,}00063$), όλα με $p<0{,}001$.
  \item Ο \textbf{PSC εξακολουθεί να ξεπερνά σημαντικά το OGD}, αλλά η διαφορά συρρικνώθηκε
        $2{,}6$ φορές μετά τη ρύθμιση του βήματος: $+0{,}00065$ πριν τη βαθμονόμηση (95\% CI
        $[+0{,}00044, +0{,}00086]$) και $+0{,}00043$ μετά, έναντι $+0{,}00149$ και
        $+0{,}00114$ με το παλιό βήμα. Το ειλικρινές συμπέρασμα παραμένει: το μείγμα ξεπερνά
        \emph{κάθε} bookmaker \emph{εκτός} του καλύτερου «sharp», όχι όλους.
  \item \textbf{Χωρίς καμία στήλη Pinnacle στο panel, το μείγμα κερδίζει τους πάντες}: με 24
        experts και οι τέσσερις αλγόριθμοι κατατάσσονται πάνω από όλους τους εναπομείναντες
        Tier-A bookmakers (OGD $0{,}99782$ έναντι VC\_BV $0{,}99987$, στους ίδιους 54.877
        αγώνες), και το OGD τους ξεπερνά σημαντικά και τους πέντε. Άρα ο περιορισμός αφορά
        \emph{μία} εξαιρετική εταιρεία, όχι τη μέθοδο --- και η Pinnacle δεν είναι καν
        προσβάσιμη σε πολλές χώρες. Με μία επιφύλαξη, βλ.\ \S\ref{sec:bestconfig}.
  \item Η \textbf{βαθμονόμηση βοηθάει σημαντικά} σε όλους τους αλγορίθμους (raw μείον
        βαθμονομημένο: $+0{,}00053$ έως $+0{,}00073$, όλα $p<0{,}001$).
\end{itemize}

\subsection{Η καλύτερη διαμόρφωση: το μείγμα φτάνει τον PSC}
\label{sec:bestconfig}

Τέσσερα πράγματα είχαν μετρηθεί χωριστά και ποτέ μαζί: το διορθωμένο βήμα του OGD, το
closing-only panel που κερδίζει το πλήρες, το Shin de-vig που κερδίζει την απλή
κανονικοποίηση στο raw log-loss, και η βαθμονόμηση. Το pipeline έτρεχε με το πλήρες panel και
την απλή κανονικοποίηση, δηλαδή άφηνε δύο από τα τέσσερα έξω. Η στοίβαξη έγινε με δύο
προφυλάξεις που το project έχει πληρώσει στο παρελθόν:

\begin{enumerate}
  \item Κάθε διαμόρφωση βαθμολογείται στους \emph{ίδιους} αγώνες --- την τομή της κάλυψης
        όλων με του PSC, 71.818 γύροι. Τα panel έχουν διαφορετική κάλυψη, και log-loss πάνω
        σε διαφορετικό σύνολο αγώνων δεν είναι συγκρίσιμο. Η \emph{εκπαίδευση} χρησιμοποιεί
        κανονικά όλο το διαθέσιμο εύρος κάθε διαμόρφωσης· μόνο η βαθμολόγηση περιορίζεται.
  \item Το να διαλέγεις διαμόρφωση κοιτώντας αυτά τα νούμερα \emph{είναι} επιλογή, οπότε όλα
        αναφέρονται δύο φορές: πλήρες δείγμα (περιγραφικά) και χρονική ουρά 25\% που η
        επιλογή δεν είδε ποτέ.
\end{enumerate}

\begin{table}[h]
\centering
\begin{tabular}{@{}lrrr@{}}
\toprule
\textbf{Διαμόρφωση} & \textbf{πλήρες} & \textbf{train 75\%} & \textbf{held-out 25\%} \\
\midrule
PSC μόνος του                        & 0,99685 & 0,99746 & 0,99502 \\
\textbf{A: closing-only + Shin}      & \textbf{0,99692} & 0,99754 & 0,99507 \\
closing-only, απλή κανονικοποίηση    & 0,99693 & 0,99756 & 0,99506 \\
baseline: πλήρες panel, απλή         & 0,99728 & 0,99801 & 0,99508 \\
πλήρες panel, Shin                   & 0,99731 & 0,99804 & 0,99512 \\
\bottomrule
\end{tabular}
\caption{Βαθμονομημένο log-loss OGD ανά διαμόρφωση, ίδιοι 71.818 αγώνες σε κάθε γραμμή.
Πηγή: \file{results/best\_configuration.csv}.}
\end{table}

Τα ζευγαρωτά τεστ δείχνουν καθαρά ποιος παράγοντας δουλεύει:

\begin{itemize}
  \item \textbf{Το panel κάνει όλη τη δουλειά.} Η διαμόρφωση A κερδίζει το baseline κατά
        $-0{,}00036$ ($p<0{,}001$)· το closing-only με απλή κανονικοποίηση κερδίζει σχεδόν
        ταυτόσημα ($-0{,}00034$, $p<0{,}001$).
  \item \textbf{Το Shin μόνο του δεν προσφέρει τίποτα μετρήσιμο} πάνω στο πλήρες panel
        ($+0{,}00004$, $p=0{,}181$) --- συνεπές με το εύρημα ότι Shin και βαθμονόμηση
        διορθώνουν το ίδιο bias, οπότε η δεύτερη το απορροφά.
  \item \textbf{Έναντι του PSC}: το baseline χάνει σημαντικά ($+0{,}00043$, $p<0{,}001$),
        ενώ η A \emph{ισοπαλεί} ($+0{,}00007$, $p=0{,}234$). Αυτή είναι η ουσιαστική αλλαγή
        στην κεντρική αφήγηση της εργασίας.
  \item Μια πέμπτη διαμόρφωση (A χωρίς τους bookmakers μίας σεζόν) βγήκε οριακά χειρότερη
        και απορρίφθηκε.
\end{itemize}

Στην held-out ουρά και οι πέντε γραμμές μαζεύονται σε λιγότερο από $0{,}0001$ μεταξύ τους.
Η ουρά επομένως \emph{επιβεβαιώνει} την ισοπαλία με τον PSC, δεν αναδεικνύει νικητή --- και
δεν θα ήταν έντιμο να παρουσιαστεί ως νίκη.

\paragraph{Η ίδια διαμόρφωση χωρίς Pinnacle: εδώ η μάθηση εξαφανίζεται.} Το πείραμα αφαίρεσης
της Pinnacle επαναλήφθηκε με την ίδια πειθαρχία (κοινοί αγώνες, train/held-out), σε δύο
panel. Στο πλήρες το αποτέλεσμα αντέχει: σε 54.877 κοινούς γύρους, και οι τέσσερις
αλγόριθμοι πάνω από όλους τους bookmakers, το OGD κερδίζει σημαντικά και τους πέντε
($-0{,}0018$ έως $-0{,}0032$, $p<0{,}001$). Στο closing-only panel όμως μένουν 12 experts και
μόνο 32.210 κοινοί αγώνες, και:

\begin{itemize}
  \item Το OGD βγαίνει \#1 ($0{,}99919$) αλλά κερδίζει σημαντικά μόνο τους 3 από τους 5
        (WHC, IWC, BWC). Ο B365C ($p=0{,}505$) και ο VC\_BVC ($p=0{,}284$) είναι ισοπαλίες,
        και στην held-out ουρά ο VC\_BVC περνάει μπροστά.
  \item Το OGD ($0{,}99919$) γίνεται πρακτικά \emph{ταυτόσημο} με τον απλό μέσο όρο
        ($0{,}99921$) --- δηλαδή η μαθημένη στάθμιση παύει να προσφέρει τίποτα.
\end{itemize}

Αυτό δεν αναιρεί το εύρημα του πλήρους panel, το εξειδικεύει: χωρίς την Pinnacle οι closing
bookmakers είναι υπερβολικά όμοιοι μεταξύ τους, οπότε δεν μένει ποιοτική διαφορά για να
μαθευτεί. Η διαμόρφωση αναφοράς είναι δηλαδή η καλύτερη εν μέρει \emph{επειδή} περιέχει τον
καλύτερο odds-setter της αγοράς.

\subsection{Τι \emph{δεν} δούλεψε (και γιατί αξίζει να αναφερθεί)}

Τέσσερις επεκτάσεις έδωσαν αρνητικό αποτέλεσμα. Τις κρατώ ρητά, γιατί η καθεμία αποκλείει
μια εύλογη υπόθεση:

\begin{itemize}
  \item \textbf{Εξειδίκευση ανά πρωτάθλημα δεν βοηθάει}: το ενιαίο μοντέλο ξεπερνά τους
        ειδικούς συνολικά ($+0{,}00058$, $p<0{,}001$) και μάλιστα σε \emph{κάθε} μία από τις
        6 λίγκες. Δύο εύλογες αιτίες: κάθε λίγκα έχει μόνο 2.231--3.800 αγώνες (πολύ λίγοι
        για να συγκλίνει ο online αλγόριθμος μόνος του), και οι ειδικοί καταλήγουν να
        εμπιστεύονται \emph{τους ίδιους} experts με το ενιαίο μοντέλο (ο PSC είναι πρώτος σε
        4 από τις 6 λίγκες), δηλαδή δεν βρίσκουν κάτι τοπικό.
  \item \textbf{Βαθμονόμηση ανά σεζόν δεν βοηθάει}, ούτε στην παραλλαγή «warm-start» που
        μεταφέρει τη θερμοκρασία της προηγούμενης σεζόν: και οι δύο κερδίζουν σημαντικά το
        raw, αλλά είναι στατιστικά \emph{αδιάκριτες} από την ενιαία global βαθμονόμηση. Ίδια
        αφήγηση bias--variance με την εξειδίκευση ανά λίγκα, σε άλλη διάσταση.
  \item \textbf{Το Fixed-Share δουλεύει, αλλά κοστίζει}: περιορίζει όντως τα διογκωμένα βάρη
        (CLC: $0{,}240 \to 0{,}064$ για $\alpha=0{,}01$· οι τέσσερις bookmakers ενεργοί μόνο
        στη σεζόν 2024/25 από $0{,}164$--$0{,}177$ σε ακριβώς $0{,}100$, δηλαδή στην
        ισοτιμία του συνόλου τους), με μικρό αλλά υπαρκτό κόστος στο συνολικό log-loss
        ($0{,}99845 \to 0{,}99885$).
  \item \textbf{Το να μαθαίνεται το ποσό Kelly online χάνει} από τον κλειστό τύπο, σε 17 από
        20 συγκρίσεις σημαντικά (και στις υπόλοιπες 3 χάνει, απλώς όχι σημαντικά). Η εξήγηση
        είναι διδακτική: ο online learner ξεκινά τυφλός ($f=0$) και μαθαίνει μόνο από
        πραγματοποιημένα κερδοφόρα/ζημιογόνα αποτελέσματα με συντηρητικό, worst-case-ασφαλή
        ρυθμό, ενώ ο κλειστός τύπος χρησιμοποιεί κατευθείαν την ήδη διαθέσιμη εκτίμηση
        πιθανότητας χωρίς κόστος εκκίνησης. Συμπέρασμα: η αξία του OCO εδώ είναι στο
        \emph{στάδιο πρόβλεψης}, όχι στο να ξαναμαθαίνει κάτι που ο κλειστός τύπος δίνει
        δωρεάν.
\end{itemize}

\subsection{Οικονομικός έλεγχος: βοηθάει το καλύτερο log-loss;}

Όχι αυτόματα --- και αυτό είναι από τα πιο ουσιαστικά ευρήματα. Καλύτερο proper scoring rule
δεν συνεπάγεται εκμεταλλεύσιμη άκρη, γιατί πρέπει να ξεπεραστεί και το περιθώριο του
bookmaker.

Ο σχεδιασμός είναι Kelly $\tfrac14$, leave-one-bookmaker-out, με βαθμονομημένες πιθανότητες,
και το \textbf{panel εκπαίδευσης ταιριάζει στη φάση του στόχου}: closing-only panel για
closing στόχο, opening-only για opening στόχο. Το δεύτερο είναι απαίτηση αιτιότητας --- καμία
closing τιμή δεν υπάρχει ακόμα όταν δημοσιεύεται μια τιμή ανοίγματος. Το πρώτο είναι επιλογή
ποιότητας, συνεπής με τη διαμόρφωση αναφοράς του \S\ref{sec:bestconfig}.

\begin{table}[h]
\centering
\small
\begin{tabular}{@{}llrrrl@{}}
\toprule
\textbf{Στόχος} & \textbf{Φάση} & \textbf{Στοιχήματα} & \textbf{Κεφάλαιο} &
\textbf{log-growth} & \textbf{$p$} \\
\midrule
WHC     & closing & 7.253  & $84{,}24\times$ & $+0{,}00061$ & $<0{,}001$ \\
IW      & opening & 7.443  & $7{,}89\times$  & $+0{,}00028$ & $0{,}008$ \\
BWC     & closing & 6.920  & $3{,}09\times$  & $+0{,}00016$ & $0{,}042$ \\
\midrule
VC\_BVC & closing & 4.561  & $1{,}58\times$  & $+0{,}00010$ & $0{,}339$ \\
WH      & opening & 5.931  & $1{,}38\times$  & $+0{,}00005$ & $0{,}545$ \\
BW      & opening & 7.841  & $1{,}09\times$  & $+0{,}00001$ & $0{,}820$ \\
B365C   & closing & 7.643  & $1{,}05\times$  & $+0{,}00001$ & $0{,}968$ \\
VC\_BV  & opening & 8.083  & $0{,}60\times$  & $-0{,}00006$ & $0{,}245$ \\
PSC     & closing & 19.497 & $0{,}54\times$  & $-0{,}00003$ & $0{,}368$ \\
B365    & opening & 10.829 & $0{,}41\times$  & $-0{,}00008$ & $0{,}101$ \\
PS      & opening & 26.675 & $0{,}26\times$  & $-0{,}00005$ & $0{,}055$ \\
\bottomrule
\end{tabular}
\caption{Μέσο log-growth ανά στοίχημα είναι η στήλη σύγκρισης· τα τελικά κεφάλαια δεν
συγκρίνονται μεταξύ γραμμών, αφού ο αριθμός στοιχημάτων διαφέρει κατά τάξη μεγέθους.}
\end{table}

\begin{itemize}
  \item \textbf{Σημαντική θετική άκρη σε 3}: WHC, IW, BWC.
  \item \textbf{Καμία σημαντική ζημιά.} Ο PSC, ο \#1 προγνώστης, βγαίνει ισοπαλία
        ($p=0{,}368$) --- το ίδιο αποτέλεσμα που δίνει και η σύγκριση πρόβλεψης στη
        διαμόρφωση αναφοράς, από ανεξάρτητο δρόμο.
  \item \textbf{Η αντιστοίχιση φάσης δεν είναι λεπτομέρεια.} Με το πλήρες panel αντί για το
        closing-only, το αποτέλεσμα χειροτερεύει σε \emph{κάθε} έναν από τους πέντε closing
        στόχους και ο PSC γίνεται σημαντική ζημιά. Προς την άλλη κατεύθυνση, μείγμα
        εκπαιδευμένο μόνο σε opening χάνει σημαντικά από \emph{κάθε} closing bookmaker
        ($-0{,}00044$ έως $-0{,}00071$, όλα $p<0{,}001$): παλιότερη πληροφορία εναντίον
        νεότερης τιμής χάνει πάντα --- το οικονομικό ανάλογο της κλίσης της κίνησης γραμμής.
\end{itemize}

Η ειλικρινής διατύπωση δεν είναι «καμία άκρη έναντι sharp, μεγάλη άκρη έναντι αδύναμων»,
αλλά: \emph{η ύπαρξη και το μέγεθος της άκρης εξαρτώνται από τον συγκεκριμένο bookmaker}. Η
προσομοίωση δεν μοντελοποιεί όρια λογαριασμού/ρευστότητας, που στην πράξη θα εμπόδιζαν τη
συστηματική εκμετάλλευση οποιασδήποτε άκρης.

\subsection{Τρία ακόμη ευρήματα με ανεξάρτητο ενδιαφέρον}

\begin{itemize}
  \item \textbf{Το προβάδισμα των closing αποδόσεων ζει σε ένα πεμπτημόριο.} Αν οι αγώνες
        ταξινομηθούν κατά το πόσο κινήθηκε η γραμμή από το opening στο closing (συνολική
        μεταβολή των τριών πιθανοτήτων) και χωριστούν σε πεμπτημόρια, η διαφορά closing
        μείον opening πάει από $-0{,}00044$ στο χαμηλότερο ($p=0{,}052$, μη σημαντική) σε
        $-0{,}01158$ στο υψηλότερο ($p<0{,}001$) --- κλίση $26\times$ μονοτονικά. Όπου η τιμή
        δεν κινήθηκε, η απόδοση κλεισίματος δεν ξέρει τίποτα παραπάνω από την αρχική. Το
        «τα closing είναι καλύτερα» δεν είναι ιδιότητα της χρονικής στιγμής, είναι ιδιότητα
        της πληροφορίας που έφτασε στο ενδιάμεσο.
  \item \textbf{Οι ισοπαλίες είναι πρακτικά απρόβλεπτες ως \emph{argmax} πρόβλεψη}: η
        καλύτερη σειρά (PSC) πιάνει μόλις $0{,}70\%$ των πραγματικών ισοπαλιών, το OGD
        $0{,}30\%$ --- παρότι οι ισοπαλίες είναι $\sim$26\% των αγώνων. Επιβεβαιώνει ποσοτικά
        ένα γνωστό φαινόμενο της ποδοσφαιρικής πρόβλεψης.
  \item \textbf{Η αγορά δεν γίνεται πιο ακριβής, γίνεται πιο ακριβή}: σε 10 σεζόν δεν
        υπάρχει σημαντική τάση σε log-loss ή ECE (όλα τα Mann-Kendall $p>0{,}37$), αλλά το
        \textbf{overround αυξάνεται σημαντικά}, από $\sim$4,4\% σε $\sim$6,0\%
        ($p=0{,}0007$) --- το αντίθετο από την αφήγηση «ώριμη, πιο ανταγωνιστική αγορά».
\end{itemize}

\begin{figure}[h]
\centering
\includegraphics[width=\textwidth]{temporal_trends.png}
\caption{Ανά σεζόν, περιορισμένο στους 4 bookmakers που υπάρχουν σε όλες τις 10 σεζόν ώστε
οι διαφορές να μην συγχέονται με αλλαγές κάλυψης. Αριστερά/δεξιά: καμία σημαντική τάση.
Μέση: το περιθώριο ανεβαίνει σταθερά.}
\end{figure}

% =====================================================================
\section{Έλεγχοι αξιοπιστίας}
\label{sec:robust}
% =====================================================================

Το μεγαλύτερο μέρος της πρόσφατης δουλειάς δεν πρόσθεσε νέα ευρήματα, αλλά έλεγξε αν τα
υπάρχοντα αντέχουν.

\paragraph{Η μέθοδος de-vig δεν αλλάζει την ιστορία.} Ξαναφτιάχτηκε ολόκληρο το dataset με τη
μέθοδο του \textbf{Shin (1992)} (μοντελοποιεί μια αναλογία «insider» στοιχηματιστών και
μετακινεί μάζα από τα outsider προς τα φαβορί, αντί να μοιράζει το περιθώριο ομοιόμορφα) και
ξανατρέχτηκε η βασική σύγκριση. Η \textbf{κατάταξη είναι ταυτόσημη} (PSC \#1, μετά
OGD$>$Hedge$>$FTRL$>$Uniform). Το Shin βελτιώνει σημαντικά τα \emph{raw} log-loss και Brier
σε όλες τις 11 οντότητες ($p<0{,}0001$) και σχεδόν υποδιπλασιάζει το ECE ($-35\%$ έως
$-57\%$· π.χ.\ OGD: $0{,}0091 \to 0{,}0043$) --- αλλά το πλεονέκτημα σε μεγάλο βαθμό
εξαφανίζεται \emph{μετά} τη βαθμονόμηση, και συμμετρικά η βαθμονόμηση παύει να είναι σημαντική
για 9 από 11 σειρές όταν έχει ήδη εφαρμοστεί Shin (έναντι 11/11 σημαντικές με proportional
normalization). Ερμηνεία: οι δύο μέθοδοι διορθώνουν το \emph{ίδιο} favorite-longshot bias με
διαφορετικό μηχανισμό --- μία ανά αγώνα, cross-sectional, η άλλη online κατά μήκος της
ιστορίας --- οπότε μόλις γίνει η μία, η άλλη γίνεται σε μεγάλο βαθμό περιττή.

\begin{figure}[h]
\centering
\includegraphics[width=0.95\textwidth]{shin_vs_basic_calibration_reliability.png}
\caption{Απόκλιση βαθμονόμησης (πραγματική $-$ προβλεπόμενη συχνότητα) ανά bin:
proportional normalization (αριστερά) έναντι Shin (δεξιά). Το εύρος συρρικνώνεται δραματικά
--- οπτική επιβεβαίωση ότι το Shin κάνει μηχανιστικά αυτό που υπόσχεται.}
\end{figure}

\paragraph{Οι αποδόσεις κλεισίματος κερδίζουν, και οι opening βλάπτουν το μείγμα.} Και για
τους 4 αλγορίθμους, ένα μείγμα \emph{μόνο} από closing bookmakers είναι σημαντικά καλύτερο
από ένα μόνο από opening ($\sim0{,}0034$--$0{,}0036$ χαμηλότερο log-loss) --- αναμενόμενο,
αφού η τιμή κλεισίματος έχει απορροφήσει την πληροφορία της κίνησης της αγοράς. Λιγότερο
αναμενόμενο: το closing-only μείγμα κερδίζει σημαντικά και το \emph{πλήρες} panel των 26,
ενώ το opening-only χάνει σημαντικά από αυτό. Δηλαδή η σειρά closing $<$ πλήρες $<$ opening
ισχύει με στατιστική βεβαιότητα, που σημαίνει ότι οι opening bookmakers είναι \emph{καθαρό
βάρος} για το μείγμα και όχι χρήσιμη διαφοροποίηση. Στη συνολική κατάταξη, το «OGD
(closing)» είναι ο καλύτερα ταξινομημένος αλγόριθμος στο \#3 συνολικά, πίσω μόνο από τον PSC
και μία σειρά μονής σεζόν. Αυτό το εύρημα είναι και ο λόγος που το closing-only panel
επιλέχθηκε τελικά ως η διαμόρφωση αναφοράς (\S\ref{sec:bestconfig}).

\paragraph{Η αιτιότητα του panel στο value betting.} Ένα σχέδιο leave-one-out δεν αρκεί να
βγάλει τη στήλη του στόχου: για opening στόχο πρέπει να φύγει \emph{όλη} η closing πλευρά,
γιατί καμία closing τιμή δεν υπάρχει ακόμα όταν δημοσιεύεται μια τιμή ανοίγματος. Αν μείνουν
μέσα, το μοντέλο στοιχηματίζει γνωρίζοντας πού θα κλείσει η αγορά, και το αποτέλεσμα είναι
τελικά κεφάλαια τάξης $10^6$--$10^9\times$, δηλαδή καθαρό τεχνούργημα. Με το panel σωστά
περιορισμένο δεν μένει μετρήσιμη άκρη σε αυτή την περίπτωση ($+0{,}00004$, $p=0{,}439$
βαθμονομημένο). Το σχέδιο τρέχει σε πλήρες grid 108 γραμμών (2 κανονικοποιήσεις $\times$ 3
τύποι panel $\times$ κάθε στόχος), ώστε κάθε συνδυασμό να μπορεί κανείς να τον δει.

\paragraph{Το βάθος ιστορίας αλλάζει συμπεράσματα.} Με 10 σεζόν, και οι τρεις αλγόριθμοι
ξεπερνούν σημαντικά την ομοιόμορφη μέση τιμή· στο μισό χρονικό παράθυρο μόνο το OGD το
πετυχαίνει. Αυτό μετρήθηκε απομονωμένα (ίδιο panel 22 λιγκών, μόνο το χρονικό παράθυρο
αλλάζει), ώστε το αποτέλεσμα να μην συγχέεται με την επέκταση σε περισσότερα πρωταθλήματα.

\paragraph{Η ανάγνωση των βαρών ενός κοιμισμένου expert.} Λόγω της αναγωγής Sleeping Experts,
ο πίνακας βαρών $W[t,i]$ συμπληρώνεται \emph{μόνο} για τους ξύπνιους experts κάθε γύρου. Άρα
η τελευταία γραμμή $W[-1]$ \emph{δεν} είναι «το τελικό βάρος κάθε expert»: είναι ακριβώς
μηδέν για όποιον δεν κάλυψε τον τελευταίο αγώνα του dataset --- 10 από τους 26 bookmakers
εδώ, μεταξύ τους ο PSC, ο καλύτερος μεμονωμένος προγνώστης. Η σωστή ποσότητα, και αυτή που
αναφέρεται παντού, είναι το βάρος που κρατούσε ο κάθε expert \textbf{την τελευταία φορά που
ήταν ξύπνιος}. Η διάκριση αφορά μόνο την ανάγνωση της κατάστασης, όχι τους ίδιους τους
αλγορίθμους: καμία μετρική ποιότητας πρόβλεψης δεν εξαρτάται από αυτήν.

% =====================================================================
\section{Γνωστοί περιορισμοί}
% =====================================================================

Αναφέρονται ρητά, δεν αποκρύπτονται:

\begin{itemize}
  \item Η συγχώνευση VC/BetVictor συνάγεται από τη μη επικαλυπτόμενη παρουσία ανά σεζόν, δεν
        είναι επιβεβαιωμένη από επίσημη πηγή --- είναι εύλογο να υπάρχουν κι άλλα μη
        εντοπισμένα rebrands, ιδίως τώρα με 22 πρωταθλήματα.
  \item Το log-loss είναι η κύρια μετρική στο μεγαλύτερο μέρος της εργασίας. Το Brier
        αναλύεται συστηματικά μόνο στη σύγκριση de-vig· το RPS υπολογίζεται παντού αλλά δεν
        αναλύεται συστηματικά πουθενά.
  \item Η προσομοίωση value betting αγνοεί όρια λογαριασμού και ρευστότητα.
  \item Ένα σχέδιο leave-one-out εισάγει look-ahead bias αν το panel δεν περιοριστεί στην
        πληροφορία που θα ήταν όντως διαθέσιμη τη στιγμή του στοιχήματος. Ο περιορισμός
        εφαρμόζεται ρητά ανά φάση αγοράς, αλλά δεν αποκλείεται να υπάρχουν άλλα, λιγότερο
        προφανή κανάλια.
  \item Οι αιτιώδεις εξηγήσεις για ορισμένα ευρήματα (π.χ.\ γιατί αυξάνεται το overround)
        είναι εύλογες, όχι αποδεδειγμένες.
\end{itemize}

% =====================================================================
\section{Κατάσταση και ερωτήματα προς συζήτηση}
% =====================================================================

Το πλήρες pipeline τρέχει καθαρά end-to-end, και το εκτενές κείμενο (θεωρία, υλοποίηση, όλα
τα ευρήματα, χάρτης αρχείων) είναι σε προχωρημένο στάδιο. Τρία ερωτήματα όπου θα ήθελα
κατεύθυνση:

\begin{enumerate}
  \item \textbf{Scope της αφήγησης.} Έχουν συσσωρευτεί $\sim$10 follow-up ευρήματα. Πόσα
        αξίζει να μπουν σε κεντρικό κεφάλαιο και πόσα σε παράρτημα, ώστε το κείμενο να έχει
        καθαρή, εστιασμένη αφήγηση αντί να διασπείρεται;
  \item \textbf{Δεύτερο άθλημα.} Θα ενίσχυε ουσιαστικά την εργασία ένα δεύτερο άθλημα
        (π.χ.\ μπάσκετ), ώστε να φανεί ότι το πλαίσιο OCO γενικεύεται πέρα από το ποδόσφαιρο;
        Το κόστος είναι κυρίως νέος αγωγός δεδομένων --- οι αλγόριθμοι δεν αλλάζουν.
  \item \textbf{Baseline προς σύγκριση.} Θα άξιζε ένα πραγματικό Bayesian Model Averaging
        baseline (συζητείται θεωρητικά, δεν έχει υλοποιηθεί), ως εμπειρικό σημείο σύγκρισης
        απέναντι στους online αλγορίθμους;
\end{enumerate}

Επιπλέον, ανοιχτά και μικρότερης κλίμακας: συστηματική αναζήτηση για κι άλλα rebrands
bookmakers στα 22 πρωταθλήματα, και διερεύνηση του \emph{γιατί} αυξάνεται το overround (θα
απαιτούσε δεδομένα εκτός football-data.co.uk).

\end{document}
